\documentclass[a4paper,11pt]{ltjsarticle}
\usepackage{amsmath,amssymb,amsthm,mathtools,hyperref,luatexja-fontspec}
\setmainjfont{HaranoAjiMincho-Regular}
\newtheorem{theorem}{定理}[section]
\newtheorem{lemma}[theorem]{補題}
\newtheorem{proposition}[theorem]{命題}
\newtheorem{remark}[theorem]{注意}
\renewcommand{\proofname}{証明}
\newcommand{\qgr}{\operatorname{qgr}}
\newcommand{\gldim}{\operatorname{gldim}}
\newcommand{\thick}{\operatorname{thick}}
\newcommand{\Perf}{\operatorname{Perf}}
\newcommand{\Db}{D^{\mathrm b}}
\title{中心三次商の一般中心値と二つの中心チャート}
\author{非可換cubic threefold研究・ソース復元用ノート}
\date{2026年11月10日・v1}
\begin{document}
\maketitle
\begin{abstract}
追跡中の\texttt{papers/pdf-only/ambient/central\_cubic\_technical.pdf}
に対応する、2026年11月1日の比較原稿第28節（pp.~43--44）から、
有限パラメータでの三次中心商の正則性を
再利用できる独立TeXとして取り出す。
使う外部のgeneric central value定理を明示し、任意の
特殊パラメータでの正則性を主張しない。
この二つの命題だけでは61ページの原PDFを置換できない。
\end{abstract}

\section{二つの中心チャートによるperfectnessの判定}
\begin{proposition}[二つの中心チャート]\label{prop:patch}
$B$をconnected graded noetherian AS-regular代数、
$W\in Z(B)$を非零の正則斉次元、$A=B/(W)$とする。
$u,v\in Z(B)_1$は$B$上正則であり、
$A$上も中心的正則列をなすと仮定する。
$A[u^{-1}]_0$と$A[v^{-1}]_0$の大域次元は有限であり、
さらに境界について
\[
\Db(\qgr A/(u,v))
=\thick\{\pi(A/(u,v))(n)\mid n\in\mathbb Z\}
\]
を仮定する。このとき
\[
\Db(\qgr A)=\thick\{\pi A(n)\mid n\in\mathbb Z\}.
\]
\end{proposition}
\begin{proof}
$A$のcoherent moduleの有限生成graded lift $N$をとり、
そのmatrix factorizationの有限presentationを用いる。
$u$と$v$を局所化した二つのaffine chartでは
有限大域次元の仮定からfactorizationがcontractibleとなる。
contractionは有限個の係数・分母から成るので分母を払え、
ある$r,s\ge1$に対して
$u^r$と$v^s$の作用がprojective singularity categoryでは零になる。

二重Koszul cone $K(u^r,v^s;N)$のcohomologyは
$(u^r,v^s)$に殺される。
\[
(u,v)^{r+s-1}\subset (u^r,v^s)
\]
なので有限filtrationを持ち、各商が$A/(u,v)$-加群となる。
境界の生成性と中心正則列のKoszul分解より、
このconeはline objectのshiftから生成されるperfect対象である。
一方、projective singularity categoryでは$u^r,v^s$の作用は
零であるため、coneは$N$の四つのshift・twistの直和に分裂する。
従って$N$自身が直和因子としてperfectになる。
各coherent $N$に適用して主張を得る。
\end{proof}

\section{固定した周囲の一般中心三次商}
\begin{theorem}[十分一般の二元三次中心係数]\label{thm:generic}
$B$をnoetherian connected graded AS-regular次元5の整域とする。
二つの次数1中心元$u,v$、三次中心元$G$を持ち、$B$が
$\mathbb C[u,v]$上自由であると仮定する。
$x,y,z$方向のfiltrationについて
\[
\operatorname{gr}_{x,y,z}B
\simeq S\otimes\mathbb C[u,v],
\]
$S$は整域、$G$の最高部分$\gamma\in S_3$は非零正則元とする。
良い中心チャートで
\[
A_{a,b,c,d}
=B/(G+au^3+bu^2v+cuv^2+dv^3)
\]
が$\mathbb C[u,v]$上自由とする。
さらに
$\qgr B/(u,v,G)$が滑らかな楕円曲線のcoherent圏と同値であるとする。
有限表示の係数全体を含む可算な部分体$k_0\subset\mathbb C$をとり、
$a,d$が$k_0(b,c)$上代数的独立であると仮定する。
このとき
\[
\gldim(\qgr A_{a,b,c,d})=3.
\]
\end{theorem}
\begin{proof}
$D_u=B[u^{-1}]_0$, $D_v=B[v^{-1}]_0$と置く。
$B$のAS-regular性と局所化から$D_u,D_v$は
noetherianで有限大域次元を持つ。
実際、$u$が中心次数1の単元で
$B[u^{-1}]\cong D_u[u,u^{-1}]$となり、
$D_u$-加群がLaurent拡大に誘導した後の
直和因子として回収できるからである。

$u$-chartに$t=v/u$を導入すると
\[
A_{a,b,c,d}[u^{-1}]_0
=D_u/(G/u^3+bt+ct^2+dt^3+a).
\]
$a$は残りの係数に対して超越的であり、
括弧内の$a$を除いた中心元は最高$x,y,z$部分が
非零の三次元なのでその冪は線形独立である。
この状況に適用されるgeneric central valueの
非可換Bertini型補題
（Rennemo--Segal--Van den Berghの中心局所化と
有限行列因子化に基づく結果）により、
このchartは有限大域次元を持つ。

同様に$v$-chartで$r=u/v$とすると
\[
A_{a,b,c,d}[v^{-1}]_0
=D_v/(G/v^3+ar^3+br^2+cr+d).
\]
$d$の超越性によってこちらも有限大域次元になる。
境界$\qgr B/(u,v,G)$の生成性と
命題\ref{prop:patch}を適用すれば
$\Db(\qgr A)=\thick\{\pi A(n)\}$。
最後に$B$のAS-regular性と中心三次正則商の
AS-Gorenstein次元4、index 2のSerre双対を使うと
大域次元は正確に3となる。
\end{proof}

\begin{remark}[適用限界と先行研究]
この命題は固定した周囲$B$について、
特定のZariski開集合全体での正則性ではなく
代数的独立性を課したvery-generalな係数の
十分条件である。各特殊係数での正則性は別問題であり、
一般混合係数の全体に滑らかな量子的境界があることも
自動的には従わない。

中心局所化・非可換generic smoothnessの外部入力として
J.~V.~Rennemo, E.~Segal, M.~Van den Bergh,
\emph{A non-commutative Bertini theorem},
\url{https://arxiv.org/abs/1705.01366}
のLemmas 3.2--3.3と関連部分を用いる。
本稿は既存研究原稿の証明をTeXで回復したものであり、
外部定理を今回改めて完全監査したものではない。
\end{remark}

\section{復元した部分と原本の関係}
原本の第28節の二つの主要命題と証明の論理をTeX化した。
原本にはその他の有限中心性、点加群、判別式、中心三次公式、
17次元モジュライと復元の議論がある。
未収録の部分を照合するまで
\texttt{central\_cubic\_technical.pdf}のGit追跡は解除しない。
\end{document}